% ============================================================================
%  F. C. Sibbern, "Logik som Tænkelære."
%  Ny Udarbeidelse, endnu som foreløbigt Aftryk for Tilhørere.
%  Kjøbenhavn: Trykt hos Directeur Jens Hostrup Schultz,
%              Kongelig og Universitets=Bogtrykker, 1827.
%  TRANSCRIPTION (Danish) — from the Royal Danish Library scan.
%
%  WHICH BOOK THIS IS — DO NOT GET THIS WRONG.
%      The title page reads *Logik som Tænkelære*. This is NOT the 1835 *Logik
%      som Tænkelære fra en intelligent Iagttagelses Standpunct og i
%      analytisk-genetisk Fremstilling*, a different and still undigitised book.
%      This is the 1827 "Ny Udarbeidelse", i.e. the SECOND EDITION of
%      *Logikens Elementer* (1822) — hence the directory name and the catalog
%      id `logikens-elementer`. Sibbern's own Fortale settles it in its first
%      sentence: "Man vil ved at sammenligne nærværende Udgave med den første
%      letteligen bemærke ...". Do not "correct" the directory or the catalog id.
%
%  SCAN: ~/bibliotek/Sibbern, Frederik/logik.pdf
%      KB / ABBYY Recognition Server, 407 PDF pages, JBIG2-encoded, ~14.6 MB.
%      KB record: "Kjøbenhavn : trykt hos Jens Hostrup Schultz, 1827. XXII, 378 s."
%      JBIG2 decodes slowly: ~1.5 s/page at 300 dpi, ~6 s at 600. `pypdf` cannot
%      extract the images at all without a `jbig2dec` binary — use `pdftoppm`.
%
%  ---------------------------------------------------------------------------
%  PAGE MAP (VERIFIED, UNIFORM — do not re-derive, do not hard-code an offset)
%
%      front matter   roman  I--XXII  ->  PDF = roman   +  3   (PDF  4-- 25)
%      body           printed 1--378  ->  PDF = printed + 25   (PDF 26--403)
%
%      Both runs uniform. No leaf was scanned twice and no leaf was dropped.
%      pagemap.py is the single source of truth.
%
%      Verified three ways:
%      (1) All 407 pages rendered at 200 dpi, top-15% strip cropped, upscaled 3x,
%          OCR'd under the Fraktur model at --psm 6. 229 pages yielded a printed
%          numeral; 205 of them give offset exactly +25. The 24 that do not are
%          all single-digit misreadings whose neighbours are +25 — chiefly a "1"
%          read as "4" in the hundreds place (PDF 128 -> "403" for 103,
%          155 -> "430" for 130, 163--165 -> "438/439/440" for 138/139/140).
%      (2) ARITHMETIC CLOSURE, which is the real proof. KB gives the extent as
%          "XXII, 378 s." Front matter PDF 4--25 is 22 pages = XXII. Body
%          PDF 26--403 is 403-26+1 = 378 = the printed extent exactly. A
%          duplicated or dropped leaf anywhere would break this.
%      (3) Endpoints read off the image: PDF 28=3, 60=35, 100=75, 200=175,
%          300=275, 400=375, 402=377; roman side PDF 10=VII, 11=VIII, 12=IX,
%          14=XI, 16=XIII, 17=XIV, 18=XV, 20=XVII, 22=XIX, 24=XXI.
%
%      DO NOT USE THE ABBYY TEXT LAYER FOR THE PAGE MAP. It was run with a Latin
%      model over Fraktur and misreads the numerals badly: a sweep over it
%      produced a spurious 20-page "+15" cluster from 5->6 confusions and a
%      "-275" cluster from contents cross-references bleeding into the first
%      line. It is a fine second witness for WORDS and useless for NUMBERS.
%
%      Structure outside the body range:
%          PDF   1     KB digitisation notice (not part of the book)
%          PDF   2-3   Kongelige Bibliotek stamp leaves
%          PDF   4     title page                     [roman I]
%          PDF   5     verso, blank                   [roman II]
%          PDF   6-7   Fortale                        [roman III--IV]
%          PDF   8-25  Oversigt (contents), ending in the Rettelser leaf
%                                                     [roman V--XXII]
%          PDF  26     printed p. 1 (numeral not printed on the page)
%          PDF 403     printed p. 378, last page of text (inside §. 176)
%          PDF 404-407 blank / back cover
%
%  ---------------------------------------------------------------------------
%  TYPE AND EMPHASIS — READ THIS BEFORE MARKING ANY EMPHASIS
%
%  TYPE: Fraktur body. Latin, foreign words, all figures, the § sign and the
%      rubric letters A.--F. are set in ANTIQUA against the Fraktur -> \textit{}.
%      This is a logic textbook and carries a great deal of inline Latin
%      (divisio, contradictio, quà, "diversus respectus tollit omnem
%      contradictionem") and single-letter variables (A, B, S, P, Non-P).
%
%  THIS BOOK DOES NOT MARK INLINE EMPHASIS BY LETTERSPACING. It marks it by
%      setting the emphasised words in SCHWABACHER against the surrounding
%      Fraktur -> \emph{}. Verified at 900 dpi on roman V by comparing the same
%      word in both faces on one page: "Erkjendelse" in §.2 and in §.4;
%      "Tænkning" in §.2 (twice) and in §.3; "Stof" in §.2 and in §.4. The
%      capitals differ unmistakably (Schwabacher E, T, S are round and open).
%      The stroke weight and ink density are the SAME in both faces, so no
%      density test separates them, and a glyph-width fit trained on the
%      pure-Fraktur Fortale separates them only partially (Schwabacher ratios
%      1.11--1.17 against Fraktur 0.92--1.04, with a contaminated middle band).
%      Schwabacher must therefore be called by eye, at 3x or better.
%
%  LETTERSPACING (Sperrsatz) -> \sperr{}, and in this book it is used for
%      DISPLAY HEADS and for the "Anm." lead-in, not for running emphasis:
%      "Logik" and "Tænkelære" on the title page, "Fortale.", "Oversigt.",
%      "Rettelser.", the Capitel labels, and "A n m." Occasional letterspaced
%      words do occur inside the Oversigt (§.170 "Definition"; the "er" of
%      §.108) and are marked.
%
%  Emphasis FOLLOWS THE PAGE, not the word: "a priori" is antiqua italic in
%      §.6 and ordinary Fraktur in §.81, and both are transcribed as printed.
%
%  EMPHASIS IN THE OVERSIGT IS DELIBERATELY NOT MARKED. Roughly a third of the
%      words in the eighteen Oversigt pages are Schwabacher. Marking them
%      reliably needs a dedicated pass at 3x over all eighteen pages; marking
%      half of them correctly is worse than marking none. What IS marked there
%      is only what is unambiguous at reading scale: \sperr{} for genuine
%      letterspacing and \textit{} for antiqua Latin. This is a known,
%      deliberate gap in the edition, not an oversight.
%
%  THE FRAKTUR DOUBLE HYPHEN is printed as "=" in compounds and is KEPT as "=":
%      Sammenhængs=Analyse, Begrebs=Analyse, Classifications=System, ideel=reel,
%      Universitets=Bogtrykker, Tænknings=Activiteter, §. 70=71.
%      Do not silently turn it into "-".
%
%  QUOTATION MARKS. The Danish low-high „...“ is the repo's default, but
%      SIBBERN'S PRINTER DOES NOT USE IT IN THE FORTALE. Both marks of the one
%      quotation on roman III are RAISED — English-style closing marks in both
%      positions — verified at 450 dpi. Transcribed there as the LaTeX ligature
%      ''...'' so that the printed form is reproduced without introducing a new
%      non-ASCII character, with a % comment at the site. Check each quotation
%      in the body against the image rather than assuming either form.
%
%  RUNNING HEADS: none in the original. The folio is an antiqua numeral centred
%      at the HEAD of the page (measured on printed p. 26). Sheet signatures sit
%      at the foot — "(b)" on roman XVII, "(b*)" on roman XIX, "(15)" on printed
%      p. 225, "(19)" on p. 26. They are not text and are not transcribed.
%
%  ---------------------------------------------------------------------------
%  STRUCTURE (from the printed Oversigt, roman V--XXII, read off the image)
%
%  There are EXACTLY TWO levels above the §: a CAPITEL level and, inside the
%  First Capitel only, a lettered A.--F. RUBRIC level. There is NO Deel and NO
%  Afdeling anywhere in the book. The Indledning (§§ 1--7) stands before the
%  First Capitel and is not part of it; the First Capitel has an Indledning of
%  its own (§§ 8--11) before rubric A.
%
%      Indledning.                                  §§   1--  7   pp.   1-- 25
%      Første Capitel.                              §§   8--100   pp.  26--245
%          Indledning.                              §§   8-- 11   pp.  26-- 29
%          A. Reel Sammenhængs=Analyse ...          §§  12-- 59   pp.  30--152
%          B. Reel Synthese, ...                    §§  60-- 67   pp. 153--178
%          C. Sammenlignings=Analyse ...            §§  68-- 74   pp. 179--207
%          D. Begrebs=Analyse.                      §§  75-- 83   pp. 208--225
%          E. Begrebers synthetiske Combination.    §§  84-- 85   pp. 226--227
%          F. Om Begreber i Almindelighed.          §§  86--100   pp. 228--245
%      Andet Capitel.                               §§ 101--128   pp. 246--311
%      Tredie Capitel.                              §§ 129--164   pp. 312--362
%      Fjerde Capitel.                              §§ 165--176   pp. 363--378
%
%  THE §§ RUN CONTINUOUSLY 1 --> 176 ACROSS THE WHOLE BOOK. They do not restart
%      at any Capitel or rubric. So a repeated \parag number is an error and a
%      hole in the range is a dropped § head. Two irregularities are in the
%      book itself, not in the transcription:
%          §§ 39 og 40 share ONE Oversigt entry and one page (101). In the print
%              "§. 39 og" stands alone on its own line above "§. 40."
%          §. 95 is SPLIT into "95 a" and "95 b" (superscript a / b), pp. 238
%              and 240. The Rettelser leaf refers to "§. 95, b", confirming it.
%      Distinct numbered Oversigt entries: 177.
%
%  A CAPITEL RANGE runs from its first §'s printed page to the page before the
%      next division's first §. The rubric ranges above were closed the same
%      way. NOTE: the Oversigt gives only the page on which each § BEGINS, so a
%      rubric HEAD's own page is not fixed by it. There are page gaps between
%      the last § of one rubric and the first § of the next (148 -> 153,
%      201 -> 208, 223 -> 226), and the rubric head stands somewhere inside the
%      gap. The skeleton below therefore does NOT pre-place the A.--F. heads;
%      each batch agent inserts \rubrik{...} where the head actually falls and
%      reports the page, so the marker range can be closed exactly.
%
%  §. HEAD FORM (measured on fourteen sampled pages: 1, 5, 6, 26, 30, 62, 101,
%      173, 246, 289, 312, 344, 363, 378 — the whole span of the book):
%      A § head is a BARE, CENTRED "§. N." standing alone on its own line. It
%      carries NO argument or rubric line beneath it: the arguments printed in
%      the Oversigt are the Oversigt's, and are not repeated at the head. The §
%      sign and the numeral are antiqua; there is a period after the sign and a
%      period after the numeral. Constant across all fourteen sampled pages.
%      Measured on printed p. 26: the head line spans x[605,697] (w=92) against
%      a measure of x[192,1097], head centre 651 against measure centre 644;
%      on p. 62 "§. 23." spans x[587,690], centre 638.
%      HENCE \parag TAKES ONE ARGUMENT (the number only) — see the macro below.
%      This is a real difference from the Nielsen books, whose \parag{n}{argument}
%      takes two. Do not copy the Nielsen signature into this book.
%
%  ANM. REMARKS. Each § is followed by one or more "Anm." remarks, numbered per
%      § and RESTARTING at 1 (Anm. 1., Anm. 2.) — so they are not checked for
%      continuity. The lead-in is LETTERSPACED FRAKTUR ("A n m.") with an
%      antiqua numeral; it is not italic and not Schwabacher.
%      THE REMARK BODIES ARE NOT SET IN SMALLER TYPE. Measured, not eyeballed:
%      on printed p. 62 the §. 22 body has a baseline pitch of 54.0 px at 300 dpi
%      against the Anm. body's 54.5 px (ratio 1.009), x-height 22 px against
%      22 px (ratio 1.000); on printed p. 1, 53.5 against 55.0 (ratio 1.028).
%      Both blocks are set to the same measure (x ~ 202-1099 against 191-1094),
%      first line indented ~87 px. Same type.
%
%  NO FOOTNOTES were found on any of the fourteen sampled pages, and Sibbern
%      cross-references inline instead ("Jfr. min Psychologie §. 59, Anm. 1"),
%      so the "Anm." remarks do the work footnotes would do. Measured with a
%      foot-region rule sweep at paper-median -15/-25/-40, gutter columns above
%      45% density masked: every wide hit came in at ink density 0.01--0.02,
%      i.e. a single-pixel scan artefact and not a rule; the full-width bands at
%      y-fraction 0.01 and 0.99 are the scanner's black page edges. No *), †) or
%      1) marks anywhere in the sample. The centred 393 px ornament at the foot
%      of printed p. 378 (y 1689--1704) is an end-of-book tailpiece, not a note
%      rule. THIS IS A SAMPLE, NOT A PROOF: if a real note turns up in a batch,
%      say so, and \renewcommand{\thefootnote} to the printed mark — NOT
%      \fnsymbol (math mode, gives U+2217) and NOT footmisc[perpage] (which
%      resets on the typeset page, not the printed one).
%
%  CAPITEL HEAD FORM, measured on printed p. 26: letterspaced Fraktur "Første
%      Capitel." centred (x[412,886]), then three centred Fraktur argument
%      lines, then a centred rule 116 px x 7 px (x[596,712]) — a short thick
%      rule, ~0.13 of the measure — then "Indledning." centred, then "§. 8."
%      The Capitel ARGUMENT is not letterspaced as a whole, but one key word in
%      it is: "Omdømmer" in the Andet Capitel argument and "Slutninger" in the
%      Tredie Capitel argument. Check each head against its page.
%
%  ---------------------------------------------------------------------------
%  ERRATA — the book's own "Rettelser", printed on roman XXII.
%
%  The Rettelser leaf is reproduced verbatim in the front matter below, since it
%  is part of the printed book. SETTLED 2026-08-26 (editor's decision): per the
%  repo's editorial stance (CLAUDE.md) the body text is transcribed AS PRINTED
%  and each erratum is logged in a "% ERRATA:" comment at its site, NOT applied.
%  Every batch agent must follow this. (This differs from
%  texts/nielsen/propaedeutiske-logik, where the errata were applied inline;
%  most of Sibbern's are substantive cross-reference corrections by the author
%  rather than compositor's slips, which is why they are left standing here.)
%  The affected body pages are:
%      24, 36, 41, 44, 49, 52, 76, 89, 99, 111, 114, 118, 170, 173, 193, 202,
%      205, 226, 233, 240, 242, 243, 289.
%  Two of them are structural and worth knowing in advance:
%      p. 240 l.18  "§. 96" should read "§. 95, b"  — confirms the 95 a/b split
%      p. 289 top   "Anm. 2 falder bort"            — an Anm. to be struck
%  The Rettelser also fix the reference "§. 62" to "§. 70=71" three times, and
%  give Sibbern's cross-reference abbreviation in LOWER CASE ("jfr.", p. 226)
%  where the body normally has "Jfr.".
%
%  ---------------------------------------------------------------------------
%  PRINTER'S DEFECTS — transcribed AS PRINTED, each flagged at its site.
%      This is a diplomatic transcription; the compositor is not corrected.
%
%      roman I    the city is spelled two ways on one page: "Kiøbenhavns" (with
%                 i) in the author's designation, "Kjøbenhavn." (with j) in the
%                 imprint. Both verified at 450 dpi. Transcribed as printed.
%      roman IV   l.11 "indbegreben" is printed with a word-space inside it,
%                 "indbe greben". SETTLED 2026-08-26: reproduced AS PRINTED, per
%                 CLAUDE.md's editorial stance. The first transcribing pass had
%                 joined the word; the join was undone. Do not re-join it — a
%                 spell-checker will want to, and it is a printed reading.
%
%  QUOTE BALANCE: 0 as of the front matter. No quote-affecting defect found yet.
%      check.py compares the raw „/“ balance against its QUOTE_DEFECTS table
%      rather than against zero, because defects CANCEL: a missing opener on one
%      page and a missing closer on another restore a raw zero that reads as an
%      all-clear. Add every new defect to that table AND to this header.
%
%  ---------------------------------------------------------------------------
%  OCR NOTES for whoever transcribes the body (the harness is untracked, so this
%  is the only place these survive).
%
%      OCR quality is GOOD under `-l Fraktur --psm 6`; the OCR can carry the
%      words and the image only has to carry structure and emphasis. The PDF's
%      own ABBYY layer is a useful SECOND WITNESS for words because it fails
%      differently (it mangles æ/ø and reads long-s as f), and is useless for
%      numbers.
%
%      THIS SCAN'S SIGNATURE ERROR: a Fraktur ";" reads as "z" — sigz, Tingz,
%          sammez. Deliberately not sed'd, because a blanket rule would eat
%          proper names and Latin; check.py flags every word-final z instead.
%          Danish native words do not end in z, so nearly every hit is real.
%      f/k/s confusion inside words: Adsfillelse->Adskillelse,
%          Modfætning->Modsætning, fættes->sættes, forfjellige->forskjellige.
%      "?" standing for a k/sk ligature: logi?->logisk, Logi?->Logik.
%      v read as o (selo->selv, objectio->objectiv); y read as v (deryed->derved).
%      dropped ø surfacing as o/s/d: gjore->gjøre, Sporgsmaal->Spørgsmaal,
%          hsre->høre, prdve->prøve, dverste->øverste.
%      F read as S at the head of a word: Sorandringer->Forandringer.
%      "§" reads as "$" ($. 41. for §. 1.).
%      Proper names break the model badly — Treschow came out as "Tre\{<ow" on
%          printed p. 315. Check every name against the image.
%      The J->I correction is a WHITELIST, not a general rule: Jeg, Ja, Jo, Jord
%          are ordinary words and "Jfr." (= jævnfør) is Sibbern's constant
%          cross-reference abbreviation.
%      Gutter noise is heavier here than on the Nielsen books: --psm 6 reads the
%          facing leaf's black edge as a stray column of | j y Y 1 ¡ » á þ at
%          BOTH margins. Ignore those columns entirely.
%      The Oversigt's OCR is NOT TRUSTWORTHY for page numbers. Both layers —
%          Fraktur tesseract AND ABBYY — agreed on a wrong printed page number in
%          five entries, all overruled by the image:
%              §  24  both read 63   image reads  68
%              §  27  both read 74   image reads  73
%              §  30  read 73 / 79   image reads  78
%              § 118  both read 280  image reads 286
%              § 149  both read 340  image reads 341
%          Confirmed shared pages (not errors): §§ 19/20 p. 59, §§ 29/30 p. 78,
%          §§ 75/76 p. 208, §§ 109/110 p. 265, §§ 135/136 p. 320,
%          §§ 151/152 p. 344, §§ 155/156 p. 347.
%
%  ---------------------------------------------------------------------------
%  SPOT-CHECK (2026-08-26, independent of the transcribing pass)
%      roman I (title page) and roman X (an ordinary Oversigt page) were
%      re-read against the scan by an agent that had not written them.
%      roman X: AGREES entry for entry. All ten entries §§ 43--52 present, none
%          duplicated, and every printed page number in the right-hand column
%          confirmed on the image (105, 106, 107, 111, 112, 113, 114, 115, 116,
%          118). Wording agrees. The "Side." column label is really printed.
%      roman I: two omissions found and CORRECTED here —
%          (i) a THIRD rule, 0.20 of the measure and ~8 px thick, stands between
%              the imprint and the date and had not been transcribed;
%          (ii) "Kjøbenhavn." in the imprint is itself letterspaced and had not
%              been marked \sperr.
%          Everything else on the page confirmed, including the two spellings of
%          the city and the three-way rule distinction.
%      The Oversigt's unmarked Schwabacher emphasis was confirmed as the known,
%      deliberate gap described above, not as an error of transcription.
%
%  STATUS
%      FRONT MATTER COMPLETE — roman I--XXII (PDF 4--25): title page, blank
%          verso, Fortale, the full Oversigt, and the Rettelser leaf, all
%          transcribed from the image.
%      BODY NOT STARTED — printed pp. 1--378. The skeleton below carries one
%          batch marker per division, eleven in all, contiguous over 1--378.
%          Batches are 12--17 printed pages, so a marker is SPLIT into
%          batch-sized markers when the batch is dispatched; splice.py matches a
%          fragment to its marker by the FIRST--LAST pair in the fragment's own
%          filename, so the split marker and the fragment name must agree
%          exactly. See the note at the head of the body skeleton.
% ============================================================================

\documentclass[12pt]{book}
\usepackage[utf8]{inputenc}
\usepackage[T1]{fontenc}
\usepackage[danish]{babel}
\usepackage[protrusion=true,expansion=true]{microtype}
\usepackage[margin=1.4in]{geometry}
\usepackage{setspace}
\usepackage{amsmath}
\usepackage{graphicx}
\usepackage{libertinus}
\usepackage{libertinust1math}
\usepackage{textalpha}   % no Greek found so far; loaded so that a later batch
                         % can type polytonic Greek directly without a preamble
                         % change mid-book.
\usepackage{fancyhdr}
% plainpages=false + pdfpagelabels: the title page, the roman front matter and
% the arabic body would otherwise all produce a PDF destination called "page.1",
% and pdfTeX drops the duplicates with a warning on every build.
% The PDF metadata is set with \hypersetup AFTER \begin{document}, NOT as
% package options here: as an option the "æ" of the title reaches hyperref while
% fontenc still has it as the T1 command \T1\ae, which hyperref cannot put in a
% PDF string ("Token not allowed in a PDF string") and silently drops.
\usepackage[colorlinks=true,linkcolor=black,urlcolor=black,
  unicode=true,plainpages=false,pdfpagelabels]{hyperref}

% A Capitel argument may carry \sperr / \textit / \\ inside it (the letterspaced
% key word "Omdømmer" in Capitel II, "Slutninger" in Capitel III), and \capitel
% puts the argument into the ToC. Those tokens are legal in the printed ToC line
% but NOT in a PDF bookmark string: hyperref warns "Token not allowed in a PDF
% string" on every build and silently drops the word. Neutralise them for the
% bookmark only, which is cleaner and more durable than hand-writing a plain
% duplicate of each argument into \texorpdfstring.
\pdfstringdefDisableCommands{%
  \def\sperr#1{#1}%
  \def\textit#1{#1}%
  \def\emph#1{#1}%
  \def\\{ }%
}

\setlength{\headheight}{15pt}
\pagestyle{fancy}
\fancyhf{}
\fancyhead[LE,RO]{\thepage}
\fancyhead[RE]{\textit{Logik som Tænkelære}}
\fancyhead[LO]{\textit{\leftmark}}
\setlength{\parindent}{1.5em}
\setlength{\parskip}{0pt}
\setstretch{1.3}

% ============================================================================
%  STRUCTURAL MACROS — use these, do not hand-roll centre blocks.
%  They exist so that every batch renders the book's recurring display forms
%  identically, and so that check.py can count them. check.py reports a
%  hand-rolled \begin{center}{\bf...} display head in \mainmatter as suspect.
%  If you need a form that has no macro here, REPORT IT and let the caller add
%  one; in particular do not hand-roll an \addcontentsline with a bare \quad in
%  it, which hyperref rejects with a warning on every build.
% ============================================================================

% Letterspacing (Sperrsatz). In THIS book letterspacing marks DISPLAY HEADS and
% the "Anm." lead-in, not running emphasis — running emphasis is Schwabacher,
% which \emph{} carries. Keep the two distinct: \sperr = letterspaced in the
% print, \emph = Schwabacher in the print.
\newcommand{\sperr}[1]{\textls[180]{#1}}

% ---------------------------------------------------------------- body heads

% A Capitel head: letterspaced Fraktur label, the argument beneath it, then the
% short thick centred rule measured on printed p. 26 (116 px x 7 px at 300 dpi,
% ~0.13 of the measure).
%   \capitel{Første Capitel.}{Om den undersøgende og behandlende ...}
% The argument is not letterspaced as a whole, but one key word inside it may be
% ("Omdømmer" in Capitel II, "Slutninger" in Capitel III) — mark that word with
% \sperr inside the argument, and check each head against its own page.
\newcommand{\capitel}[2]{%
  \par\vspace{1.3\baselineskip}%
  \begin{center}
    {\large\sperr{#1}\par}\vspace{0.45\baselineskip}
    {#2\par}
  \end{center}
  \vspace{0.35\baselineskip}
  \begin{center}\rule{0.13\linewidth}{1.4pt}\end{center}
  \vspace{0.45\baselineskip}
  \phantomsection
  \addcontentsline{toc}{chapter}{\texorpdfstring{#1\ #2}{#1 #2}}%
  \markboth{#1}{#1}%
}

% A lettered rubric head, A.--F., inside the First Capitel. The letter is
% printed in ANTIQUA on its own centred line above the rubric heading; the
% heading may run to several centred lines, separated with \\.
%   \rubrik{A.}{Reel Sammenhængs=Analyse\\eller\\Reelle Gjenstandes Analyse ...}
% The rubric heads are NOT pre-placed in the skeleton: the Oversigt gives only
% the page on which each § begins, so a rubric head's own page falls somewhere
% in the gap before the next §. Insert it where it actually stands and report
% the page.
\newcommand{\rubrik}[2]{%
  \par\vspace{1.1\baselineskip}%
  \begin{center}
    {\large\textit{#1}\par}\vspace{0.35\baselineskip}
    {#2\par}
  \end{center}
  \vspace{0.5\baselineskip}
  \phantomsection
  \addcontentsline{toc}{section}{\texorpdfstring{#1\ #2}{#1 #2}}%
}

% A centred group head standing alone — used for "Indledning.", both the
% book's own (before the First Capitel) and the First Capitel's internal one.
\newcommand{\gruppe}[1]{%
  \par\vspace{0.9\baselineskip}%
  \begin{center}{#1\par}\end{center}
  \vspace{0.4\baselineskip}
  \phantomsection
  \addcontentsline{toc}{section}{#1}%
}

% A § head. ONE ARGUMENT — the number only.
% MEASURED on fourteen pages spanning the whole book: a § head is a bare,
% centred "§. N." on its own line, antiqua sign and numeral, with NO argument
% line beneath it. The arguments printed in the Oversigt are the Oversigt's and
% are not repeated at the head. Do not give this macro a second argument.
%   \parag{12}      ->   §. 12.
% Irregular numbers occur and are passed through as text:
%   \parag{39 og\\ \S.~40}   (two lines, as printed)
%   \parag{95 \textsuperscript{a}}
\newcommand{\parag}[1]{%
  \par\vspace{0.9\baselineskip}%
  \begin{center}{\textit{\S.\,#1.}\par}\end{center}
  \vspace{0.35\baselineskip}
  \noindent\ignorespaces
}

% An "Anm." remark lead-in. Letterspaced Fraktur "A n m." with an antiqua
% numeral; the remark BODY is set in the SAME type size as the § body (measured
% — see the header), so no \small here. Anm. numbers restart at 1 in each §.
%   \anm{1}   ->   Anm. 1.        \anm{}   ->   Anm.
\newcommand{\anm}[1]{%
  \par\vspace{0.5\baselineskip}%
  \noindent\sperr{Anm.}\ifx&#1&\else~\textit{#1}.\fi\enspace\ignorespaces
}

% ------------------------------------------------------- Oversigt (contents)
%  The Oversigt is reproduced as a FACSIMILE of the printed contents, not as a
%  structural apparatus; this edition's own page numbers come from the generated
%  \tableofcontents that follows it. check.py's hand-rolled-head test runs only
%  over \mainmatter for exactly this reason.

\newcommand{\ovhead}[1]{%
  \begin{center}{\Large\sperr{#1}}\\[1.2ex]\rule{0.30\linewidth}{0.4pt}\end{center}
  \vspace{0.6\baselineskip}
}

% The right-aligned "Side." column label, printed once at the head of every
% Oversigt page.
\newcommand{\ovside}{\par\nopagebreak\hfill{\footnotesize Side.}\par\nobreak\vspace{0.2\baselineskip}}

\newcommand{\ovgroup}[1]{%
  \par\vspace{0.7\baselineskip}\begin{center}{#1\par}\end{center}
  \vspace{0.3\baselineskip}}

\newcommand{\ovcap}[2]{%
  \par\vspace{1.0\baselineskip}%
  \begin{center}{\sperr{#1}\par}\vspace{0.3\baselineskip}{#2\par}\end{center}
  \vspace{0.4\baselineskip}}

\newcommand{\ovrubrik}[2]{%
  \par\vspace{0.9\baselineskip}%
  \begin{center}{\textit{#1}\par}\vspace{0.25\baselineskip}{#2\par}\end{center}
  \vspace{0.4\baselineskip}}

\newcommand{\ovnote}[1]{%
  \par\vspace{0.5\baselineskip}\begin{center}\begin{minipage}{0.86\linewidth}
  \centering #1\end{minipage}\end{center}\vspace{0.4\baselineskip}}

% One § entry, leader-dotted to its printed page number.
%   \ovline{12}{Det givne Mangfoldige maa først analyseres ...}{30}
\newcommand{\ovline}[3]{%
  \par\vspace{0.25\baselineskip}%
  \noindent\hangindent=2.8em\hangafter=1
  \textit{\S.\,#1.}\quad #2%
  \nobreak\ \leaders\hbox to 0.55em{\hss.\hss}\hfill\ \textit{#3}\par}

\hyphenation{Sam-men-hængs-Ana-lyse Sam-men-lig-nings-Ana-lyse
  Grund-fore-stil-lin-ger Er-kjen-del-sen Frem-gangs-maa-der
  Til-freds-stil-lel-se Be-grebs-be-stem-mel-ser Clas-si-fi-ca-tions-Sy-stem}

\begin{document}
\hypersetup{
  pdftitle={Logik som Tænkelære (1827)},
  pdfauthor={F. C. Sibbern},
  pdfsubject={Transcription from the Royal Danish Library scan}
}
\frontmatter
\pagestyle{empty}

% =====================================================================
%  F. C. Sibbern, *Logik som Tænkelære* (1827)
%  FRONT MATTER, roman I--IV  =  PDF pp. 4--7
%      roman I    title page
%      roman II   verso, blank
%      roman III  Fortale, first part
%      roman IV   Fortale, conclusion, dated and signed
%
%  Transcribed from the image (KB scan, 300 dpi, with 450--900 dpi zooms
%  on the title-page imprint and the quotation in the Fortale).
%
%  TYPOGRAPHIC CONVENTIONS USED HERE -- macros NOT yet in transcription.tex.
%  The caller must add them (or substitute):
%      \sperr{...}   letterspaced (Sperrsatz) words. Used on this book for
%                    DISPLAY heads only -- "Logik", "Tænkelære", "Fortale.",
%                    "Oversigt.", the Capitel heads. \emph{} is an acceptable
%                    stand-in if a separate macro is not wanted.
%      \textit{...}  Latin/foreign set in antiqua against the Fraktur.
%  This book's INLINE emphasis is NOT letterspacing: it is SCHWABACHER set
%  against the Fraktur (see the header of oversigt.texfrag). No inline
%  emphasis occurs in the Fortale.
%
%  QUOTATION MARKS. Sibbern's printer does NOT use the Danish low-high
%  „...“ here. Both marks of the one quotation on roman III are RAISED
%  -- English-style closing marks in both positions -- verified at 450 dpi.
%  Transcribed as the LaTeX ligature '' ... '' so that no new non-ASCII character is introduced
%  and check.py's „/“ balance stays at zero. See the % comment at the site.
%
%  The Fraktur double hyphen is kept as "=" in compounds, per the book's
%  conventions (Universitets=Bogtrykker).
% =====================================================================


% --- roman I ---
% Title page, PDF 4. Centred display Fraktur throughout. Two rules:
%   (a) a small lozenge/leaf ornament rule under the author block;
%   (b) a thick-over-thin double rule above the imprint.
% "Logik" and "Tænkelære" are letterspaced; "Tænkelære" is the largest
% type on the page (117 px against 90 px at 450 dpi). "Kjøbenhavn." in
% the imprint is letterspaced too — found on the 2026-08-26 spot-check,
% missed by the first pass.
% THREE rules stand on this page, and they are distinct ornaments:
%   0.42 of measure  lozenge with hairline extensions, under the author block
%   1.00 of measure  thick (~8 px) over thin (~3 px), ~7 px apart, above the
%                    imprint
%   0.20 of measure  a plain ~8 px rule between the imprint and the date
% All three measured at 450 dpi over a 1363 px measure.
% NOTE: the page spells the city two ways -- "Kiøbenhavns" (with i) in the
% author's designation, "Kjøbenhavn." (with j) in the imprint. Both verified
% at 450 dpi; transcribed as printed, not normalised.
\begin{center}
  {\Huge\sperr{Logik}}\\[2.2ex]
  som\\[2.2ex]
  {\Huge\sperr{Tænkelære}}\\[3.5ex]
  af\\[1.8ex]
  {\large Dr. Frederik Christian Sibbern,}\\[0.4ex]
  {\footnotesize Professor i Philosophien ved Kiøbenhavns Universitet.}\\[3ex]
  % ornament rule (lozenge with hairline extensions)
  \rule{0.42\linewidth}{0.4pt}\\[3.5ex]
  {\large Ny Udarbeidelse,}\\[0.6ex]
  endnu som foreløbigt Aftryk for Tilhørere.\\[4ex]
  % thick-over-thin double rule
  \rule{\linewidth}{1.4pt}\\[-1.2ex]
  \rule{\linewidth}{0.4pt}\\[4ex]
  {\large\sperr{Kjøbenhavn.}}\\[0.8ex]
  {\large Trykt hos Directeur Jens Hostrup Schultz,}\\[0.3ex]
  {\footnotesize Kongelig og Universitets=Bogtrykker.}\\[2.2ex]
  % A THIRD rule, between the imprint and the date: 279 px = 0.20 of the
  % 1363 px measure, ~8 px thick, at 450 dpi. Added on the 2026-08-26
  % spot-check; the first pass missed it. Distinct from the 0.42-measure
  % lozenge ornament above and from the full-measure thick-over-thin pair.
  \rule{0.20\linewidth}{1.0pt}\\[2.2ex]
  1827.
\end{center}


% --- roman II ---
% Verso of the title page. Blank; no printed numeral, no signature.


% --- roman III ---
% PDF 6. No printed numeral on this page: the display head stands instead.
% "Fortale." is letterspaced, with a lozenge ornament rule beneath it.
% The first word opens with a large two-line ornamental initial M.
\begin{center}
  {\Large\sperr{Fortale.}}\\[1.6ex]
  \rule{0.30\linewidth}{0.4pt}
\end{center}

Man vil ved at sammenligne nærværende Udgave med den første letteligen
bemærke, at, om end Grundanskuelsen og Fremgangsmaaden i det Hele er bleven
den, den var, saa er nærværende Udgave dog saa forskjellig fra hiin, at
Forholdet imellem begge, naar det skulde tilkjendegives nøiagtigt, maatte
udtrykkes ved at kalde denne en ny Udarbeidelse, som da jo ogsaa den første
Udgave blot var udgivet som ''Manuskript for Tilhørere,'' fordi
% Both quotation marks here are RAISED in the print (English-style closing
% marks in both positions), not the Danish low-high „...“.
% Verified at 450 dpi. Set as '' ... '' so the printed form
% is reproduced without a new non-ASCII character.
Forf. ønskede at kunne gjennemarbeide den endnu engang. Men ved den
Omarbeidelse, den nu er undergaaet, er Bogen dog endnu ikke, hvad Forfatteren
ønskede, at den skulde være, eller hvad han troer, at kunne, under fortsat
Eftertænkning, gjøre den til; hvorfor han ikke ønsker den betragtet som gjort
færdig fra hans Haand. Han har derfor --- ligesom han strax, da han lod
Trykningen begyndes, besluttede kun at lade et, for en Bog som denne, ikke
stort Oplag foranstalte deraf --- ogsaa endnu ikke villet lade den komme for
Lyset anderledes, end som foreløbigt Aftryk, for sine Tilhøreres Skyld, uagtet
han denne Gang vil lade den komme i Boghandelen.

Forfatteren bliver ved at benytte den Leilighed, hans Stilling giver ham, til
at lade sit Skrivt undergaae den Prøve, et Skrivt kommer til at blive
underkastet, naar der holdes gjentagne Forelæsninger derover, en Prøve, der
torde være en af de solideste af alle; især da den vil kunne befordre den, som
tænkende Venner ville kunne lade det vederfares. Hvad der især staaer tilbage,
er for det første, i Henseende til Realiteten, en mere tilfredsstillende
Opklaring af nogle enkelte Partier, som endnu ei ere fremstillede
fyldestgjørende, dernæst ogsaa en nærmere Bestemmelse i Henseende til
adskillige Problemer, som Anordningen frembyder, f. Ex. den rette Plads for
Læren om et Begrebs Determinationer og om Definitioner og
% --- roman IV ---
% PDF 7. No printed numeral on this page either.
Inddelinger; men endeligen staaer, hvad Bogen heelt igjennem angaaer, endnu
tilbage en nøiere Omsorg for at give den de Fortrin i Henseende til Foredraget
og det Enkeltes Fremstilling, eller hvad man kalder Formen, hvilke den endnu
maa savne. Man vil bemærke, at Forfatteren af denne Bearbeidelse af Logiken
har forsøgt, hvad man ved historiske Værker kalder: paa ny at udarbeide dem
efter Kilderne. Idet han vilde forsøge at fremstille Logiken saaledes, at den
kunde fortjene Navn af Tænkelære --- et Navn, som Prof. Twesten i Kiel har
været resolut nok til, at nægte den sædvanlige Logik, sin egen indbe greben
% PRINTER'S DEFECT, REPRODUCED AS PRINTED: "indbegreben" carries a word-space
% inside it, "indbe greben" (roman IV, l. 11). Compositor's spacing, not a
% variant spelling. The first pass joined the word and logged the join; per
% CLAUDE.md's editorial stance the compositor is not corrected silently, so the
% split is now reproduced and logged here instead. (Editor's decision,
% 2026-08-26.) Do not "fix" this.
--- og søgte at tage Alt saa meget mueligt fra første Haand, ved umiddelbar
Iagttagelse af Tænkningen, især i dens videnskabelige Stræben, uden at lade sin
Gang i denne Henseende bestemmes ved den Gang, Andre havde gaaet, gav Sagen ham
saameget at gjøre, at Formen derover for det første kom til at staae tilbage;
thi til at kunne tilfredsstille de Fordringer, man desto heller gjør til sig
selv i Henseende til denne, jo mere hiin er bleven En kjær og magtpaaliggende,
udfordres, at man allerede maa være saa raadig over sit Stof, at man kan bevæge
sig deri med Lethed. Havde det ikke for mine Forelæsningers Skyld været
nødvendigt, at lade den nærværende Udarbeidelse trykke strax, saa havde jeg
saare gjærne ladet den blive beroende hos mig selv, til jeg kunde have ladet
den undergaae den Regeneration, som jeg, dersom det forundes mig at opleve et
nyt Oplags Fornødenhed, haaber da at kunne lade den finde.

\begin{center}
  Kjøbenhavn den 20\textsuperscript{de} Juli 1827.
\end{center}
\begin{flushright}
  Sibbern.
\end{flushright}

\cleardoublepage
\pagestyle{fancy}

% =====================================================================
%  F. C. Sibbern, *Logik som Tænkelære* (1827)
%  OVERSIGT (contents), roman V--XXII  =  PDF pp. 8--25
%  Ends with the book's Rettelser (errata) leaf, roman XXII.
%
%  Transcribed FROM THE IMAGE, page by page (playbook §6), with the
%  Fraktur-tesseract layer and the PDF's own ABBYY layer used only as two
%  independent second witnesses. Every printed page number was read off the
%  image; the two OCR layers agreed on a wrong number in five places
%  (§§ 24, 27, 30, 118, 149) and both were overruled by the image.
%
%  MACROS ASSUMED -- none of these is in transcription.tex yet. Add them, or
%  substitute; the fragment is otherwise plain text.
%      \ovhead{...}                 the "Oversigt." display head (letterspaced)
%      \ovside                      the right-aligned "Side." column label,
%                                   printed once at the head of every page
%      \ovgroup{...}                a centred group head ("Indledning.")
%      \ovcap{head}{argument}       a Capitel head with its argument block
%      \ovrubrik{A.}{lines}         a lettered rubric (A.--F.) with its heading
%      \ovnote{...}                 a parenthetical editorial note, centred and
%                                   set between entries
%      \ovline{n}{argument}{page}   one § entry, leader-dotted to the page no.
%      \sperr{...}                  letterspaced words
%      \textit{...}                 Latin/foreign set in antiqua
%
%  ------------------------------------------------------------------
%  EMPHASIS -- READ THIS BEFORE ADDING ANY \emph TO THIS FILE.
%
%  This book does NOT mark inline emphasis by letterspacing. It marks it by
%  setting the emphasised words in SCHWABACHER against the surrounding
%  Fraktur. Verified at 900 dpi on roman V by comparing the same word in both
%  faces on one page: "Erkjendelse" in §.2 and in §.4; "Tænkning" in §.2
%  (twice) and in §.3; "Stof" in §.2 and in §.4. The capitals differ
%  unmistakably (Schwabacher E, T, S are round and open); the stroke weight
%  and the ink density are the SAME in both faces, so no density test
%  separates them, and a glyph-width fit trained on the pure-Fraktur Fortale
%  separates them only partially (Schwabacher ratios 1.11--1.17 against
%  Fraktur 0.92--1.04, with a contaminated middle band).
%
%  Roughly a third of the words in the Oversigt are Schwabacher. Marking them
%  reliably needs a dedicated pass at 3x or better over all eighteen pages;
%  that was beyond this batch's budget, and marking half of them correctly is
%  worse than marking none. SO: NO INLINE EMPHASIS IS MARKED BELOW. What is
%  marked is only what is unambiguous at reading scale --
%      \sperr{} for genuine letterspacing (display heads; §.170 "Definition";
%               the "er" of §.108), and
%      \textit{} for Latin set in antiqua.
%  Emphasis follows the page, not the word: "a priori" is antiqua italic in
%  §.6 but ordinary Fraktur in §.81, and both are transcribed as printed.
%  ------------------------------------------------------------------
%
%  The Fraktur double hyphen is kept as "=" in compounds
%  (Sammenhængs=Analyse, Classifications=System, ideel=reel, §. 70=71).
%
%  Sheet signatures at the foot -- "(b)" on roman XVII, "(b*)" on roman XIX --
%  are not text and are not transcribed.
% =====================================================================


% --- roman V ---
\ovhead{Oversigt.}

\ovgroup{Indledning.}
\ovside
\ovline{1}{Begreb om Logik. Dens Forskjel fra den psychologiske Lære om
  Forstanden. Benævnelser. Forskjel imellem naturlig og videnskabelig Logik:
  speculativ og beskrivende Logik.}{1}
\ovline{2}{Om den umiddelbare Erkjendelse, som gaaer i Forveien for Tænkning,
  og giver den dens Stof eller Materie. Begreb om Tænkning i almindelig
  Forstand}{5}
\ovline{3}{Forskjel imellem Tænkningens Form og dens Gehalt. Logik som Lære om
  den videnskabelige Tænkning. Hvorvidt den er en blot formal Videnskab. Om den
  skal være et Videnskabens Organon, eller blot dens Kanon}{6}
\ovline{4}{Begreb om Tænkning i videnskabelig Forstand. Den derved vundne
  Erkjendelse er en middelbar. Modsætning imellem Anskuelse og Tænkning, og det
  i to Henseender. Nærmere Betragtning af Tænkningens forskjellige Stof}{10}
\ovline{5}{Hvilke Hovedpuncter der ere at bemærke, naar der er Spørsmaal om
  Tænkningens almindeligste Formaal i det Hele. Herved om den dobbelte Betydning
  af Ordet Sandhed, og om den forskjellige Væren eller Væsenhed (Realitet), der
  kan tillægges en Gjenstand}{13}

% --- roman VI ---
\ovside
\ovline{6}{Hvorledes Tænkningens Gang og Maade bestemmes paa den ene Side ved
  visse Grundforestillinger og Grundforudsætninger \textit{a priori}, som ligge
  til Grund for den}{20}
\ovline{7}{Paa den anden Side ved Beskaffenheden af det Givne, hvorpaa den
  gaaer ud}{25}

\ovcap{Første Capitel.}{Om den undersøgende og behandlende Tænknings
  forskjellige Fremgangsmaader, og den Erkjendelse, som derved opnaaes.}

\ovgroup{Indledning.}
\ovline{8}{Det umiddelbart Givne er enten noget Enkelt eller noget
  Mangfoldigt}{26}
\ovline{9}{Hvorvidt det Enkelte kan give Tænkningen noget Stof}{26}
\ovline{10}{Hvorledes det Mangfoldiges Betragtning kan føre til en tredobbelt
  Undersøgelse og Behandling}{27}
\ovline{11}{De hertil svarende Tænkningens Fremgangsmaader. At disse deels
  kunne have reelle Anskuelser, deels Begreber til Gjenstande}{29}

\ovrubrik{A.}{Reel Sammenhængs=Analyse\\
  eller\\
  Reelle Gjenstandes Analyse ifølge Grundbegrebet om reel Sammenhæng.}
\ovline{12}{Det givne Mangfoldige maa først analyseres. Forskjellig Brug af
  dette Ord}{30}
\ovline{13}{Det første Skridt er her Anskuelsens Fuldendelse. Hvorpaa den
  beroer}{32}

% --- roman VII ---
\ovside
\ovline{14}{Derved erholdes tillige en orienterende Almeenforestilling. Om og
  hvorvidt denne er et Begreb}{34}
\ovline{15}{Det næste Skridt er, formedelst en intellectuel Analyse, at komme
  til Indsigt i Tingens Sammenhæng, m. m., og derved at komme til Begreb om den.
  Nærmere om denne nye Erkjendelses Forhold til Anskuelsen og Iagttagelsen, og
  dens eiendommelige Beskaffenhed. Hvad Art af Sammenhæng her betragtes. Det
  Aprioriske, som her ligger til Grund}{37}
\ovline{16}{Om Begrebets og Begribningens Udvikling, og hvorledes den kan føre
  til et heelt reelt Begrebssystem}{43}
\ovline{17}{Dobbelt Art af Begriben og Begreb: Sammenhængsbegreb og
  Realbegreb. Nærmere om begges Forhold til hinanden og Beskaffenhed}{44}

\ovnote{(Nu følger en nærmere Udvikling af Maaden, hvorpaa vi komme til den
  intellectuelle Erkjendelse og af dennes egen Beskaffenhed, samt
  Beskaffenheden af Det i Tingene, hvorpaa den gaaer ud.)}

\ovline{18}{Trende Tænknings=Activiteter, som indtræde under den betragtede
  Fremgangsmaade. Hvorledes de gribe organiskt i hinanden}{55}
\ovline{19}{Forskjel imellem dem og selve den Erkjendelse, hvortil de føre, og
  som formedelst dem fremstilles}{59}
\ovline{20}{Til denne Erkjendelse komme vi paa en middelbar Maade}{59}
\ovline{21}{Men en umiddelbar intellectuel Skuen ligger dog derved til Grund,
  og hiin Erkjendelse beroer paa den ene Side herpaa}{60}
\ovline{22}{Men paa den anden Side beroer den paa en Tænkningens
  Tilfredsstillelse}{61}
\ovline{23}{Denne maa finde Sted baade i formel og i materiel Henseende.
  Forskjel imellem at gjøre sig et Begreb om Noget, og at begribe det. Hvorledes
  Begreb kan erholdes paa en mere eller mindre middelbar eller umiddelbar Maade.
  Herved om Hypotheser}{62}

% --- roman VIII ---
\ovside
\ovline{24}{Hvad der giver den intellectuelle Erkjenden Charakteren af en
  Erkjenden. Om flere Slags Nødvendighed, og den derpaa beroende forskjellige
  Art af Erkjendelse. Bemærkning angaaende Ordet: Erkjendelse}{68}
\ovline{25}{Erkjendelsen betragtet som grundet. Forskjel imellem to Slags
  Grunde}{71}
\ovline{26}{Hvori en Tings Grund egentligen ligger}{72}
\ovline{27}{Hvorledes vi, under den betragtede Fremgangsmaade, komme til den
  intellectuelle Erkjenden ifølge Raisonnement og tilmed paa en discursiv
  Maade}{73}
\ovline{28}{Tænkningen, idet den søger Tilfredsstillelse, søger det
  Forstandsmæssige i Tingene, hvilket nu og kan kaldes noget Intellectuelt.
  Herved om den forskjellige Betydning af Ordet Begreb, da det tages baade i
  objectiv og i subjectiv Forstand}{75}
\ovline{29}{Det Intellectuelle i Tingen kan og kaldes det Ideelle i den}{78}
\ovline{30}{Dette Intellectuelle og Ideelle er at betragte som en egen
  Gjenstand, en Tænkningsgjenstand, for sig}{78}
\ovline{31}{Det hæves frem som saadant formedelst en Abstraction, og Begrebet
  træder derved frem som abstract, under den her beskrevne Tænkningens Gang er
  det tillige abstraheret}{81}
\ovline{32}{Forskjel imellem Betragtning \textit{in abstracto} og
  \textit{in concreto}. Intuition}{82}
\ovline{33}{Forholdet imellem det Intellectuelle i Tingen og dens Udvortes.
  Hiint er det i sig selv Primitive og Constitutive, dette det i sig selv
  Secundære og Consecutive i Tingen. Lignende Forhold imellem Tingens
  Realbegreb og dens Sammenhængsbegreb}{85}

% --- roman IX ---
\ovside
\ovline{34}{Der indtræder da nu en Omvending, ifølge hvilken hiint i sig selv
  Primitive nu bliver Det, hvorfra der gaaes ud}{87}
\ovline{35}{Idet Gjenstanden herved i og for Anskuelsen paa en eftergjørende
  Maade construeres, eller genetiskt fremstilles, danner der sig et Tankeheelt,
  hvilket --- idet Gjenstandens Indhold derved Led for Led udtydes, og det
  tilsvarende Begrebssystem Led for Led deduceres --- selv kan siges at blive
  construeret. Nærmere Bemærkning over, hvad et Begrebs Construction vil
  sige}{88}
\ovline{36}{Hvorledes det Intellectuelle og til Grund Liggende i Tingen træder
  frem for Betragtningen som Tingens Natur og Væsen. Bemærkninger over disse
  Udtryk. Om Experimenteringen}{91}
\ovline{37}{Hvorledes Begrebet som det til Grund Liggende medfører visse
  Begrebsbestemmelser; for det første \textit{sensu activo}. Hvorledes enhver af
  dem kan træde frem for sig, og antage Charakteren af et
  Beskaffenhedsbegreb}{97}
\ovline{38}{Om Begrebsbestemmelser \textit{sensu passivo}, og Begrebets
  Bestemmelse ved dem}{100}
% §§ 39 and 40 share one entry and one page number; "§. 39 og" stands alone on
% its own line above "§. 40." in the print.
\ovline{39 og\\ §.~40}{Dobbelt Svar paa det Spørsmaal: Hvori Begrebet
  egentligen bestaaer. Paa den ene Side bestaaer det deri, at der indsees og
  begribes; paa den anden Side bestaaer det i det Indbegreb af
  Begrebsbestemmelser, hvilke erkjendes som tilsammenhørende, og sees i deres
  Eenhed}{101}
% DOUBTFUL READING: the "af" before "Begrebsbestemmelser" sits in the gutter at
% the extreme left of roman IX and is barely inked; tesseract read "1", ABBYY a
% comma. "af" is the only reading that construes ("det Indbegreb af
% Begrebsbestemmelser"); the rejected alternative was a bare comma with no
% preposition.
\ovline{41}{Hvad et Begrebs Indhold er}{103}
\ovline{42}{Om det i intellectuel Henseende Væsentlige, og det endnu desuden
  Tiltrædende}{104}

% --- roman X ---
\ovside
\ovline{43}{\textit{Essentialia constitutiva} og
  \textit{essentialia consecutiva}}{105}

\ovnote{(Nu følger en nærmere Betragtning af den discursive og discuterende
  Tænkning (§. 27), som udfordres for at komme til Begreber, og af den
  Begrebsudvikling, hvortil den fører).}

\ovline{44}{Hvorledes Begribningen, som en Indtrængen i det Forstandsmæssige,
  der er i Tingen selv, hviler i et dobbelt Støtte= og Udgangspunct}{106}
\ovline{45}{Hvorledes under den Undersøgelse, som fører til Begrebet, visse
  Forstandsfordringer danne sig og søge Fyldestgjørelse. Noget indenfra
  Kommende mødes med det udenfra Givne}{107}
\ovline{46}{Hvorledes vi herunder orientere os i Gjenstandens Betragtning, og
  til hvilke Adskillelser og Indsigtens Udvidelser denne Orientering kan
  føre}{111}
\ovline{47}{Den orienterende Adskillen af Gjenstandens Hovedsider og
  Hovedpartier}{112}
\ovline{48}{Den partielle Begriben og Begrebsdannelse, som deraf vil kunne
  blive en Følge}{113}
\ovline{49}{Om Modsætningen hertil: det altomfattende Begreb, eller
  Totalbegrebet}{114}
\ovline{50}{Hvorledes der atter, hvad dette sidste angaaer, er at skjelne
  imellem Grundbegrebet og det fuldelige Totalbegreb eller det totale
  Specialbegreb. Den gradvise Begriben}{115}
\ovline{51}{I Henseende til Grundbegrebet skjelne vi igjen imellem det første
  og nærmeste Grundbegreb og det totale Grundbegreb. Hvorledes de herved
  adskilte Hensyn kunne staae hinanden mere eller mindre nær}{116}
\ovline{52}{Hvorledes Begrebet kan gaae ud over det Givne. Først ved at
  uddanne sig til Normalbegreb}{118}

% --- roman XI ---
\ovside
\ovline{53}{Dernæst ved at overføres paa ganske andre Gjenstande, end de
  givne, og tjene til disses Opklaring og Forstaaelse}{120}
\ovline{54}{Hvorledes det Foregaaende fører til at skjelne imellem en
  Gjenstands Grundnatur med den tilsvarende Grundtypus, og dennes særdeles
  Modificationer i Gjenstanden. Forskjel imellem det relativ Væsentlige og det
  relativ Tilfældige i Gjenstanden. Hvorledes hiin Skjelnen imellem Grundnatur
  og Modificationer viser sig i sin Gyldighed paa flere Maader. Dobbelt Art af
  Modificerende. Hvorledes den i det Specielle gaaende intellectuelle
  Indtrængen i een Gjenstands hele Beskaffenhed, for at udgrunde dennes
  Hvorledes og Hvorfor, kan føre ind i en stor Kreds af Undersøgelser, idet den
  i samme Grad trænger dybere ind}{122}
\ovline{55}{Grundbegrebet træder frem som et Begreb for sig, formedelst en
  Abstraction, som vi kalde den blot orienterende eller relative. Om den
  objective Gyldighed af en saadan Abstraction. --- Den eensidige orienterende
  Abstraction; den blot eensidige Abstraction; de fixe Abstractioner eller
  Betragtningerne \textit{in abstracto} efter fixe Synspuncter}{132}
\ovline{56}{Grundbegrebet har sit eget Indhold, men dette er mindre end
  Specialbegrebets. Mere og mindre fundamentelle Begreber; mere og mindre
  sammensatte}{137}
\ovline{57}{Grundbegrebet som det Primitive, og som Udgangspunctet for en
  logisk Deduction og et logiskt Begrebssystems Dannelse}{140}
\ovline{58}{Hvorledes et Grundbegreb i Almindelighed har Charakteren af
  generelt. Om Specialbegrebet kan dette ei gjelde. Om en virkelig Fælledshed
  for flere kan siges at være væsentlig for et Grundbegreb. To Arter af saadan
  Fælledshed; høiere og lavere Begreber}{141}

% --- roman XII ---
\ovside
\ovline{59}{Hvorvidt den hele Classe af de, paa den hidtil forfulgte Vei
  erholdte, intellectuelle Erkjendelser eller Begreber gaaer ud paa det mere
  Værende, og derfor er af høiere Gyldighed, end de reelle Anskuelser}{148}

\ovrubrik{B.}{Reel Synthese,\\
  eller\\
  reelle Anskuelsers eller Gjenstandes Frembringelse ved en selvstændig
  synthetisk Fremgangsmaade.}
\ovline{60}{Selvstændig synthetisk Fremgangsmaade. Flere Slags Exempler
  derpaa}{153}
\ovline{61}{Dertil udkræves: visse Elementer (Enkeltheder), hvoraf --- et
  Schema, hvori --- en reel Activitet, hvorved --- og en Norm (et Begreb),
  hvorefter der construeres. Constructionsbegreber; hvad de have tilfælleds med
  de forhen betragtede Begreber}{157}
\ovline{62}{Hvorledes vi komme til denne Classe af Begreber. At Gjørligheden
  af Det, et saadant Begreb gaaer ud paa, maa godtgjøres, hvor det ei indlyser
  af sig selv}{160}
\ovline{63}{Hvorvidt med disse Begreber en sand Begriben er forbunden; hvad
  der giver dem Charakteren af egentlige Sammenhængsbegreber. --- Hvorledes der
  i den construerede Anskuelse eller Gjenstand kan vise sig at ligge Mere, end
  vi ved Constructionen have nedlagt i den; hvorledes dette Mere er at udfinde
  og erkjende}{163}

% --- roman XIII ---
\ovside
\ovline{64}{Hvorledes og hvorvidt der ad denne Vei danne sig Realbegreber, og
  hvilke Forskjelle der i denne Henseende ere at bemærke}{168}
\ovline{65}{Hvorvidt her de i den construerede Anskuelse eller Gjenstand
  indbefattede Enkeltheder blive Gjenstande for særegne Begreber}{172}
\ovline{66}{Hvorledes under den regelmæssige synthetiske Fremgangsmaades Gang
  høiere og mere fundamentelle Begreber danne sig tidligere, end de lavere og
  mere specielle. Forskjellen imellem Væsentligt og Tilfældigt anvendt her.
  Hvori Begrebets høiere Frihed viser sig her}{175}
\ovline{67}{Om de her betragtede Begrebers objective Gyldighed eller
  Realitet}{178}

\ovrubrik{C.}{Sammenlignings=Analyse\\
  eller\\
  Gjenstandes og Begrebers Analyse og Ordning ifølge Grundbegrebet om Lighed
  og Sammenhæng formedelst Lighed.}
\ovline{68}{Hvorpaa denne Tænkningens Virksomhed i Almindelighed gaaer ud;
  hvad dertil udfordres}{179}
\ovline{69}{Hvorledes den heri indtrædende og til Grund liggende Erkjendelse
  af det Lige i ellers ulige Gjenstande viser sig som en Erkjendelse af ideel
  og intellectuel Art. Flere Beskaffenheder ved den. Hvad der til Fremhævelsen
  af hiint Lige i det ellers Ulige udfordres}{181}
\ovline{70}{Hvorvidt den sig herved dannende Erkjendelse er et Begreb}{186}
\ovline{71}{Om og hvorvidt dette bestaaer i en Indsigt. Definition paa
  Almeenbegreb. Hvorledes ogsaa Erkjendelsen af det Ulige, betragtet som det
  Charakteristiske antager en begrebsmæssig Charakteer. Hvorvidt
  Almeenbegreberne gaae ud paa det mere Værende, og derfor ere af høiere
  Betydenhed, end de individuelle Anskuelser. Om de kunne hæve sig til en
  høiere Frihed ved at gaae ud over det Givne}{188}

% --- roman XIV ---
\ovside
\ovline{72}{Hvad der udfordres, for at Lighedernes og Ulighedernes Erkjendelse
  skal kunne føre til en bestemt Ordnen. Hvorledes deres større eller mindre
  Væsentlighed bestemmes og bedømmes. Egentlige Slægtsbegreber. Hvori deres
  Indhold bestaaer}{194}
\ovline{73}{Logisk Over= og Underordning. Slægter og Arter.
  Classifications=System. Forskjel imellem Classification i synoptisk og i
  heuretisk Henseende; naturlig og konstig Classification}{198}
\ovline{74}{Hvorledes Tænkningen ad denne Vei endvidere kan føre til en
  Sammenhængs= og Realbegriben, idet Rækker eller Rækkesystemer dannes.
  Sammenhængen er her at kalde ideel=reel. Denne kan undertiden vise sig ved
  Siden af den egentligen reelle. Til et Rækkesystem kunne vi komme deels ad
  Iagttagelsens, deels ad den selvstændige Productions Vei}{201}

\ovrubrik{D.}{Begrebs=Analyse.}
\ovline{75}{Spørsmaalet er nu, hvorvidt de forhen beskrevne Fremgangsmaader
  kunne anvendes paa Begreber}{208}
\ovline{76}{Simpel intellectuel Analyse}{208}
\ovline{77}{Begrebsanalyse blot som Forfølgelse af et Begrebs
  Conseqvents}{209}
\ovline{78}{Den mere egentlige Begrebsanalyse}{211}
\ovline{79}{Hvorledes denne for det første blot kan gaae ud paa Begrebet
  saaledes, som det er givet i den almindelige Forestillingsmaade, og udtaler
  sig i Sprogbrugen}{212}

% --- roman XV ---
\ovside
\ovline{80}{Hvorledes denne Analyse i sin Forfølgelse, naar den skal føre til
  Begrebets grundige Bestemmelse, maa gaae tilbage til en reel
  Sammenhængs=Analyse, eller til en Sammenlignings=Analyse}{214}
\ovline{81}{Den egentlige rene Begrebsanalyse. Den gaaer ud paa de almindelige
  Grundbegreber a priori, og de speciellere Begreber, hvori de træde frem under
  særdeles Modificationer. Saadanne ere deels de rene almindelige
  Forstandsbegreber, hvilke vi og kunne kalde de rene Reflectionsbegreber,
  hvortil alle de almindelige Forholdsbegreber henhøre, deels Ideerne,
  forsaavidt disse træde frem i Begrebers Form}{215}
% "a priori" here is set in FRAKTUR, not in the antiqua italic used for it in
% §. 6. Transcribed as printed; not italicised.
\ovline{82}{At ogsaa denne Art af Analyse maa, for at Erkjendelsen kan
  fuldende sig, føre til en Deduction og synthetisk Fremgangsmaade; men at
  dette dog ei ophæver den logiske Realitet af denne Art af
  Begrebsanalyse}{220}
\ovline{83}{Den hidtil betragtede Begrebsanalyse kan kaldes den objective
  eller reelle Begrebsanalyse. Om en derfra aldeles forskjellig subjectiv,
  eller blot logisk Begrebsanalyse, der ogsaa kan kaldes den blot
  formelle}{223}

\ovrubrik{E.}{Begrebers synthetiske Combination.}
\ovline{84}{Hvori den bestaaer; hvorledes den maa skee}{226}
\ovline{85}{Om de derved dannede Begrebers Realitet}{227}

% --- roman XVI ---
\ovside
\ovrubrik{F.}{Om Begreber i Almindelighed.\\
  (Resultater af det Foregaaende.)}
\ovline{86}{Almindelig Definition paa Begreb}{228}
\ovline{87}{Tre Slags Begreb}{229}
\ovline{88}{Hvad de indbyrdes have Fælleds og Forskjelligt}{230}
\ovline{89}{Om det Dialectiske ved de to af disse tre Arter, nemlig ved
  Almeen= og ved Realbegreberne}{231}
\ovline{90}{Almindelig Udsigt over, hvorledes vi komme til enhver af de tre
  Arter af Begreber}{233}
\ovline{91}{Fire almindelige Dannelsesmaader}{234}
\ovline{92}{De tre Hovedactiviteter, som vise sig i den undersøgende
  Fremgangsmaade. Af Combination i Særdeleshed gives der fire Arter}{234}
\ovline{93}{At Begrebsdannelsen dog egentligen beroer, deels paa intellectuel
  Skuen, deels paa conseqvent Tankeforfølgelse og Tænknings Tilfredsstillelse.
  --- Hvorvidt Begreber ere discursive Erkjendelser}{235}
\ovline{94}{Begrebers Inddeling efter deres Kilde i Begreber
  \textit{a posteriori} og Begreber \textit{a priori}. Hovedclasserne af de
  sidste. Absolut og relativ Aprioritet}{236}
\ovline{95 \textsuperscript{a}}{Forskjel imellem Begreb og dets Gjenstande,
  Begreb og reel Anskuelse}{238}
\ovline{95 \textsuperscript{b}}{Forskjel imellem Begrebet selv og de
  Modificationer, hvorunder det træder frem}{240}
\ovline{96}{Begrebets Indhold og Omfang}{241}
\ovline{97}{Grundbegreber og Specialbegreber, høiere og lavere}{242}
\ovline{98}{Disses indbyrdes Forhold i Henseende til Indhold og Omfang}{243}
\ovline{99}{De forskjellige Begrebssystemer}{244}
\ovline{100}{At Begreber overalt træde frem under Form af Omdømmer}{245}

% --- roman XVII ---
\ovcap{Andet Capitel.}{Om Begrebers og overhovedet Erkjendelsers Fremstilling
  ved Hjelp af \sperr{Omdømmer} og disses Natur og Beskaffenhed i
  Almindelighed.}
\ovside
\ovline{101}{Omdømmet som væsentlig Fremtrædelsesform for Begreb og al ifølge
  Tænkning sig dannende Erkjendelse. Forskjel imellem egentlige Omdømmer og
  Sætninger med Omdømmets Form}{246}
\ovline{102}{Hvad der i et Omdømme er at adskille. Dobbelt Betydning af
  Ordet: Prædicat. Definition paa Omdømme. Kryptiske Sætninger}{249}
\ovline{103}{Hvorledes den Erkjendelse, som i et Omdømme udtaler sig,
  forholder sig til dette selv, og hvor den er at søge i samme}{251}
\ovline{104}{Omdømmet nærmere betragtet: det er 1) en Position og kun
  een;}{253}
\ovline{105}{2) en Henførelse af en Forestilling til en anden; 3) en
  Bestemmelse af en Forestilling ved Hjelp af den anden; 4) en Forbindelse af
  en Forestilling med en anden: to Maader, hvorpaa Omdømmet kan
  betragtes}{254}
\ovline{106}{Et Omdømme er endeligen 5) en saadan Sammensættelse af Tvende til
  Eet, at de dog deri blive staaende som Tvende}{256}
\ovline{107}{Hvad det vil sige, at Subjectet siges at være Prædicatet.
  Forskjellig Kraft af dette Udtryk}{258}
\ovline{108}{Dobbelt Betydning af det logiske \sperr{er} i Henseende til den
  reale Nexus imellem Subject og Prædicat}{263}
\ovline{109}{Den nærmere Betragtning af Omdømmet gaaer ud fra at betragte det
  som Resultat af en Tænkning}{265}
\ovline{110}{Hvorledes da Modsætningen imellem et Ja og et Nei træder frem.
  Prædicat i vidtløftig og i indskrænket Betydning}{265}

% --- roman XVIII ---
\ovside
\ovline{111}{Hvorledes Omdømmer overhovedet kunne inddeles i fire
  Hoved=Henseender}{269}
\ovline{112}{Almindelig Bemærkning angaaende de formelle Inddelinger. Hvad der
  ved enhver af dem i det Følgende er at iagttage}{270}
\ovline{113}{Inddelingen i Henseende til Qvaliteten}{271}
\ovline{114}{Hvorledes vi komme til flere formelle Inddelinger}{276}
\ovline{115}{Inddeling i Henseende til Qvantiteten}{277}
\ovline{116}{De Hovedarter af Omdømmer, der vise sig, naar begge de anførte
  Inddelinger forbindes med hinanden}{280}
\ovline{117}{Inddeling i Henseende til Modaliteten. Qvalitetsforskjellenes
  Indflydelse paa denne Inddeling. Almindelige Slutningsregler, som her ere at
  bemærke}{281}
\ovline{118}{Inddeling i Henseende til Relationen. Nærmere om det hypothetiske
  Omdømmes Beskaffenhed. To Maader, paa hvilke samme kan betragtes}{286}
\ovline{119}{Det disjunctive Omdømme}{291}
\ovline{120}{Hvorledes Relationsinddelingen lader os see to ganske
  forskjellige Arter af Forhold, som ved et Omdømme ere at ponere}{292}
\ovline{121}{Om enkelte og sammensatte Sætninger}{293}
\ovline{122}{Om flere Omdømmers indbyrdes Overeensstemmelse og
  Ikke=Overeensstemmelse, og de logiske Inddelinger af Omdømmer, som i denne
  Henseende ere at gjøre}{295}
\ovline{123}{Omdømmers forskjellige Beskaffenhed med Hensyn til den reelle
  Nexus imellem Subject og Prædicat. Omdømmers forskjellige Gehalt}{298}
\ovline{124}{Forskjel imellem de egentlige dømmende og de ikke egentligen
  dømmende Sætninger}{300}

% --- roman XIX ---
\ovside
\ovline{125}{Til de sidste høre de blot begrebsopstillende og de blot
  beskrivende Sætninger}{301}
\ovline{126}{Fremdeles de identiske Sætninger. Dobbelt Betydning af dette
  Udtryk}{302}
\ovline{127}{Om de egentlige dømmende Sætninger: disse ere enten analytiske
  eller synthetiske. Flere Maader, hvorpaa analytiske Sætninger kunne danne
  sig}{304}
\ovline{128}{Overgang til Læren om Slutninger}{311}

\ovcap{Tredie Capitel.}{Om Maaden, hvorpaa Erkjendelser kunne vise sig som
  grundede og gyldige, i Almindelighed, og om \sperr{Slutninger} i
  Særdeleshed.}

\ovline{129}{Den første og øverste af de logiske Grundsætninger, som i denne
  Henseende ere at bemærke: Identitetsprincipet}{312}
\ovline{130}{Contradictionsprincipet}{314}
\ovline{131}{Udelukkelsesprincipet}{317}
\ovline{132}{At disse Grundsætninger ere blot formelle. Ethvert Omdømme maa
  have en materiel Grund}{317}
\ovline{133}{Et Omdømmes Sandhed kan enten være givet umiddelbart eller
  middelbart}{318}
\ovline{134}{I første Tilfælde er det enten givet \textit{de facto} eller som
  Axiom. To Slags Axiomer. Om ikke enhver Sætning, hvorved et Factum udsiges,
  kan siges at danne sig ved Slutning}{319}
\ovline{135}{Hvad der forstaaes ved at bevise: to Slags Beviis}{320}
\ovline{136}{Begreb om Slutning; hvad deri er at bemærke}{320}
\ovline{137}{Almindelig Regel for alle Slutninger}{324}

% --- roman XX ---
\ovside
\ovline{138}{Hvad der hører Logiken til, i Læren om Slutninger at
  oplyse}{325}
\ovline{139}{Hvorledes vi i Almindelighed kunne gaae frem, for at udfinde og
  bestemme de forskjellige Slags Slutninger i formel Henseende}{326}
\ovline{140}{Om de umiddelbare Slutninger i Almindelighed}{329}
\ovline{141}{Om Identitets=, Underordnings= og Modsætningsslutningen i
  Særdeleshed}{330}
\ovline{142}{Begreb om Conversionsslutning}{331}
\ovline{143}{Dens forskjellige Arter}{332}
\ovline{144}{Hvorledes hypothetiske og disjunctive Sætninger kunne blive
  Gjenstande for Conversionsslutning}{333}
\ovline{145}{Om den middelbare Slutning i Almindelighed}{334}
\ovline{146}{Dens første Art, den kategoriske Syllogismus, i Almindelighed
  betragtet}{335}
\ovline{147}{Den naturligste Fremgangsmaade i samme: den første Figur}{336}
\ovline{148}{Hvorledes i alt fire Figurer lade sig tænke}{339}
\ovline{149}{Den anden Figur}{341}
\ovline{150}{Den tredie Figur}{342}
\ovline{151}{Den fjerde Figur}{344}
\ovline{152}{Almindelige Regler for alle kategoriske Syllogismer}{344}
\ovline{153}{Den middelbare Slutnings anden Hovedart. Den almindelige
  hypothetiske Slutning}{345}
\ovline{154}{Speciel Art af hypothetisk Slutning}{346}
\ovline{155}{Den middelbare Slutnings tredie Hovedart: den disjunctive
  Slutning. Apagogiskt Beviis}{347}
\ovline{156}{Dilemma}{347}
\ovline{157}{Sammenkjædede Slutninger. Sorites}{348}

% --- roman XXI ---
\ovside
\ovline{158}{Enthymema. Epichirema}{349}
\ovline{159}{Slutninger betragtede fra den reale Side. Udsigt over de
  forskjellige Arter af Beviser}{350}
\ovline{160}{Inductions= og Analogieslutning}{352}
\ovline{161}{Forklaringsslutninger. Beviis af Indicier}{354}
\ovline{162}{Demonstrationerne}{356}
\ovline{163}{De begrebsforfølgende Beviser. De ere enten simple analytiske
  Beviser, eller debatterende og discuterende Argumentationer, eller
  Deductioner}{357}
\ovline{164}{Adskillige Feil, som ved Slutninger maae undgaaes, i Henseende
  til Det, hvorfra, og Det, hvorpaa der gaaes ud}{360}

\ovcap{Fjerde Capitel.}{Om videnskabelig Fremstilling i Almindelighed og
  Beskrivelse, Definition og Inddeling i Særdeleshed.}

\ovline{165}{Begreb om Videnskab i Almindelighed}{363}
\ovline{166}{Forskjellig Methode i den videnskabelige Fremstilling}{364}
\ovline{167}{Begrebers Forfølgelse i to Retninger, med Hensyn til det dobbelte
  Slags Sammenhæng}{365}
\ovline{168}{Hvad der ligger til Grund for al Videnskab, som systematisk}{366}
\ovline{169}{Om de to Slags Beskrivelse}{367}
\ovline{170}{Begreb om \sperr{Definition}. Hvorledes den er forskjellig fra
  Beskrivelse}{368}
\ovline{171}{Almindelige Regler for Definitioner}{370}
\ovline{172}{Flere Arter af Definition. Hvad der hører til at udfinde
  Definitioner. Hvorledes vi ledes hen til dem ad flere Veie}{371}

% --- roman XXII ---
\ovside
\ovline{173}{Inddeling i Almindelighed. Den logiske Inddeling i Særdeleshed.
  Hvad der ved samme er at bemærke}{374}
\ovline{174}{Regler for samme. --- Inddeling efter Leddenes Antal}{376}
\ovline{175}{Reel Inddeling}{377}
\ovline{176}{Distinction imellem en Gjenstands forskjellige Sider.
  Bemærkninger angaaende flere Slags Distinction}{378}

% Lozenge ornament rule, then the errata leaf. The Rettelser are the LAST
% printed matter of the front matter; PDF 26 begins printed p. 1. There is no
% second errata leaf at the end of the volume (PDF 404--407 are blank / cover).
\begin{center}
  \rule{0.34\linewidth}{0.4pt}\\[2.5ex]
  {\large\sperr{Rettelser.}}
\end{center}

% Printed as a three-column list: the page number, the line number, and the
% correction; the repeated words "Pag." and "Lin." are replaced by a dash in
% every line after the first. Set here as plain lines, since only the wording
% matters for the edition.
\begin{flushleft}
Pag. 24 Lin. 16, istedetfor Lit. C læs: B, §. 62.\\
--- 36 --- 2, istedetfor §. 62 læs: §. 70=71.\\
--- 41 --- 4, istedetfor B læs: C.\\
--- 44 --- 22, istedetfor §. 62 læs: 70.\\
--- 49 --- 3, efter desl. sluttes Parenthesen.\\
--- 52 --- 8 fra neden, den blotte formale Begriben læs: den blotte
Sammenhængsbegriben.\\
--- 76 --- 12, hvilket læs: hvilke.\\
--- 89 --- 19 og 20, den læs: det.\\
--- 99 --- 18, isted. §. 62 læs: §. 70=71.\\
--- 111 --- 6 fra neden, eller læs: deels til en.\\
--- 114 --- 16 fra oven, kunne de Ord og skeer styrkeviis falde bort.\\
--- 118 --- 12, efter kunne sættes: undertiden.\\
--- 170 --- 9, isted. §. 73 læs: §. 74.\\
--- 173 --- 9, isted. §. 73 læs: §. 74, Anm. 5.\\
--- 193 --- 23, isted. §. 73 læs: §. 74.\\
--- 202 --- 4, isted. og læs: af.\\
--- 205 --- 6, isted. §. 70 læs: §. 71.\\
--- 226 --- 13 fra neden, efter 57 sættes: jfr. §. 77, Anm.\\
--- 233 --- 6 fra neden, efter Sammenhængsbegreber sættes: saavelsom
Realbegreber.\\
--- 240 --- 18, maae istedetfor §. 96 sættes: §. 95, \textsuperscript{b}.\\
--- 242 --- 11 fra neden, §. 72 læs: §. 73.\\
--- 243 --- 16, 83 læs: 84.\\
--- 289 øverste Linie, Anm. 2 falder bort.
\end{flushleft}
% The Rettelser confirm that the book has Anm. remarks numbered per § (p. 173
% "§. 74, Anm. 5"; p. 289 "Anm. 2") and that Sibbern's cross-reference
% abbreviation is "jfr." in lower case here (p. 226), not "Jfr.".

\cleardoublepage
\tableofcontents

\mainmatter    % sets arabic numbering itself — do not add \pagenumbering here

% ============================================================================
%  BODY — printed pp. 1--378 (PDF 26--403).
%
%  One batch marker per division, derived from the
%  printed Oversigt. The eleven ranges below are contiguous and cover 1--378
%  exactly, so check.py's gap test is meaningful from the first batch on.
%
%  HOW TO USE A MARKER. Batches are 12--17 printed pages, so a division marker
%  is normally SPLIT into batch-sized markers before dispatch: the single marker
%  for pp. 30--152 becomes a stack of markers of the same form reading
%  30--44, 45--59, 60--74 and so on, until the division is covered.
%  (No literal marker text is written out in this comment, deliberately: a
%  second copy of a marker string anywhere in the file makes splice.py refuse to
%  run, because it insists on finding exactly one copy of the marker it wants.)
%  splice.py matches a fragment to its marker by the FIRST--LAST pair parsed out
%  of the fragment's own FILENAME (pp30-44.texfrag), so the split marker and the
%  fragment name must agree exactly, character for character. Never edit
%  transcription.tex from inside a batch agent.
%
%  HEADS. The Capitel heads and the two "Indledning." heads are pre-placed
%  below, because the Oversigt fixes their pages (a Capitel begins on the page
%  of its first §, and printed p. 26 was verified on the image). The A.--F.
%  rubric heads are NOT pre-placed: the Oversigt gives only the page on which
%  each § begins, and there is a gap between the last § of one rubric and the
%  first § of the next (148->153, 201->208, 223->226), so each rubric head
%  stands somewhere inside its gap. The batch agent inserts \rubrik{...} where
%  the head actually falls, reports the page, and the range is closed then.
%  If a head turns out to stand EARLIER than the marker range assumed, say so
%  and the marker boundary moves.
% ============================================================================

\gruppe{Indledning.}
% §§ 1--7.  Printed p. 1 carries no folio (the numeral is not printed on it).
% Division marker pp. 1--25 split for dispatch, 2026-08-26.
% § starts inside this range (from the Oversigt): §.1 p.1, §.2 p.5, §.3 p.6,
% §.4 p.10, §.5 p.13, §.6 p.20, §.7 p.25.
% [text to be added: pp. 1--13]
% [text to be added: pp. 14--25]


\capitel{Første Capitel.}{Om den undersøgende og behandlende Tænknings
forskjellige Fremgangsmaader, og den Erkjendelse, som derved opnaaes.}
% Head verified on the image, printed p. 26 (PDF 51): letterspaced Fraktur
% label centred x[412,886], three centred Fraktur argument lines, then a centred
% rule 116 px x 7 px (x[596,712]) at 300 dpi, then "Indledning.", then "§. 8."

\gruppe{Indledning.}
% §§ 8--11, the First Capitel's own Indledning, before rubric A.
% [text to be added: pp. 26--29]

% --- rubric A. --- §§ 12--59.  Insert here, at its true page:
%     \rubrik{A.}{Reel Sammenhængs=Analyse\\
%       eller\\
%       Reelle Gjenstandes Analyse ifølge Grundbegrebet om reel Sammenhæng.}
% Division marker pp. 30--152 split for dispatch, 2026-08-26. Only the first
% batch is split off so far; split the remainder as batches are dispatched.
% § starts in the first batch: §.12 p.30, §.13 p.32, §.14 p.34, §.15 p.37,
% §.16 p.43, §.17 p.44.
% [text to be added: pp. 30--44]
% [text to be added: pp. 45--152]

% --- rubric B. --- §§ 60--67.  The head falls between p. 148 (§. 59) and
%     p. 153 (§. 60); insert it where it stands and report the page:
%     \rubrik{B.}{Reel Synthese,\\
%       eller\\
%       reelle Anskuelsers eller Gjenstandes Frembringelse ved en selvstændig
%       synthetisk Fremgangsmaade.}
% [text to be added: pp. 153--178]

% --- rubric C. --- §§ 68--74:
%     \rubrik{C.}{Sammenlignings=Analyse\\
%       eller\\
%       Gjenstandes og Begrebers Analyse og Ordning ifølge Grundbegrebet om
%       Lighed og Sammenhæng formedelst Lighed.}
% [text to be added: pp. 179--207]

% --- rubric D. --- §§ 75--83.  The head falls between p. 201 (§. 74) and
%     p. 208 (§. 75):
%     \rubrik{D.}{Begrebs=Analyse.}
% [text to be added: pp. 208--225]

% --- rubric E. --- §§ 84--85.  The head falls between p. 223 (§. 83) and
%     p. 226 (§. 84):
%     \rubrik{E.}{Begrebers synthetiske Combination.}
% [text to be added: pp. 226--227]

% --- rubric F. --- §§ 86--100:
%     \rubrik{F.}{Om Begreber i Almindelighed.\\
%       (Resultater af det Foregaaende.)}
% [text to be added: pp. 228--245]


\capitel{Andet Capitel.}{Om Begrebers og overhovedet Erkjendelsers Fremstilling
ved Hjelp af \sperr{Omdømmer} og disses Natur og Beskaffenhed i Almindelighed.}
% §§ 101--128.  "Omdømmer" is letterspaced in the Oversigt's form of this
% argument; VERIFY it against the printed head on p. 246 and, if the head
% differs from the Oversigt, record BOTH readings rather than normalising them
% (repo CLAUDE.md, editorial stance).
% [text to be added: pp. 246--311]


\capitel{Tredie Capitel.}{Om Maaden, hvorpaa Erkjendelser kunne vise sig som
grundede og gyldige, i Almindelighed, og om \sperr{Slutninger} i Særdeleshed.}
% §§ 129--164.  Same caution: "Slutninger" letterspaced per the Oversigt;
% verify against the printed head on p. 312.
% [text to be added: pp. 312--362]


\capitel{Fjerde Capitel.}{Om videnskabelig Fremstilling i Almindelighed og
Beskrivelse, Definition og Inddeling i Særdeleshed.}
% §§ 165--176.  Printed p. 378 is the last page of text and ends inside §. 176;
% it carries a centred 393 px tailpiece ornament at the foot (y 1689--1704 at
% 300 dpi), which is an end-of-book ornament and not a footnote rule.
% [text to be added: pp. 363--378]


\backmatter
\end{document}
